\section{Related Work}
\label{sec:related}

\paragraph{LLM-guided invariant and lemma generation.}
The use of large language models to generate candidate loop invariants has attracted
significant recent attention. CIll~\cite{cill2026} is the closest prior work to
VGuide: it uses CTI (counterexample trace) context to guide an LLM in generating
invariants for hardware model checking. Like VGuide, CIll exploits the CE trace
rather than static source analysis, and its LLM oracle is subject to external
validation before use. The primary difference is the verification domain---hardware
model checking of RTL circuits versus software CEGAR over C programs---and the
abstraction mechanism: CIll targets transition-system invariants for symbolic model
checkers, while VGuide targets predicate-abstraction refinement in a CEGAR loop.
The hybrid LLM+SMT framework of~\cite{loopinvariant2025} addresses software loop
invariant synthesis by alternating between LLM candidate generation and SMT
validation, achieving strong results on standard benchmarks. That work focuses on
standalone invariant synthesis rather than integration into the CEGAR refinement
loop; its validation step is analogous to VGuide's L3 tier, but it operates outside
a running verifier. Quokka~\cite{quokka2025} accelerates program verification by
using LLMs to synthesize auxiliary invariants that are then checked by a
portfolio of solvers; it reports gains comparable to VGuide on overlapping benchmarks
and shares the insight that LLM-generated invariants must be independently
certified. VGuide differs from Quokka in its CE-triggered query strategy and
its direct integration into CPAChecker's CEGAR loop. Work on synthesizing
inductive invariants for distributed protocols~\cite{ic3llm2026} applies the LLM
oracle principle to IC3-based model checking of distributed systems---the oracle
architecture is closely related to VGuide, with the LLM proposing invariant
candidates that IC3 checks; the domain difference (distributed protocols and TLA+
vs.\ C programs and predicate abstraction) makes the technical challenges distinct.
Large Lemma Miners~\cite{largelemma2025} explores whether LLMs can generate inductive
lemmas for hardware verification, finding that while LLMs can identify structurally
similar lemmas, they struggle with deep inductive reasoning, motivating the
validate-before-inject approach that VGuide also adopts.
Not All Invariants Are Equal~\cite{notallequal2026} takes an orthogonal direction:
rather than using LLMs at inference time to generate predicates, it curates training
data to improve the quality of a small language model specialized for invariant
synthesis. This training-time approach complements VGuide's inference-time oracle and
could in principle be used to fine-tune a local model for VGuide's predicate oracle.

\paragraph{Neuro-symbolic and agentic verification.}
At a broader scale, neuro-symbolic proof generation~\cite{neurosymbolic2026} uses
LLM guidance to drive interactive theorem provers (Lean, Coq) for systems software
verification. The integration philosophy resembles VGuide---LLM provides hints,
formal tool validates---but the target is interactive proof construction rather than
predicate abstraction, and the feedback cycle operates at the level of proof tactics
rather than counterexample traces. Agentic model checking~\cite{agenticmc2026}
proposes wrapping a model checker in an LLM agent loop: the agent iteratively
selects configurations, interprets results, and mutates the checking strategy. VGuide
is deliberately lighter-weight: it makes a single structured oracle call per task and
injects the result into a standard CEGAR loop, without autonomous agent logic or
iterative replanning. This design choice reduces the attack surface for LLM
misbehavior and keeps the latency overhead minimal. The use of LLMs to extract
problem structure for SAT local search~\cite{satllm2025} shares the philosophy of
using learned models to provide high-quality starting points for combinatorial
solvers; VGuide can be understood as the CEGAR analogue of that idea, where the
``starting point'' is a set of predicates at loop heads rather than a variable
assignment for local search.

\paragraph{Foundation: CEGAR and predicate abstraction.}
VGuide is built directly on the CEGAR paradigm~\cite{cegar2003} and predicate
abstraction~\cite{predabstraction1997}. Craig-interpolation-based
refinement~\cite{interpolation2003} is the standard predicate discovery mechanism
that VGuide supplements at loop heads. The SLAM project~\cite{slam2002} demonstrated
the practical utility of predicate abstraction for industrial software, establishing
the paradigm that CPAChecker~\cite{cpachecker2011} implements and that VGuide extends.
VGuide is implemented within CPAChecker and participates in the SV-COMP
ecosystem~\cite{svcomp2022}. The soundness arguments in \Cref{sec:design} rely
directly on the theoretical foundations of these works: VGuide does not modify the
abstraction--model-check--refine cycle, and its predicate injection step respects
the monotonicity of predicate abstraction that underpins the CEGAR soundness proof.
